% Background & methodology: everything the rest of the report assumes.

This section is the vocabulary for the rest of the report. It explains how
the game works, how the bots decide, what distinguishes each bot version,
and how we score things --- so that later sections can say what happened
without stopping to define terms.

\subsection{The game, in one paragraph}
\label{sec:bg-game}

Players compete for territory on a random graph. Every node can hold troops;
nodes with factories produce troops each tick, and powerplants boost
adjacent factories. Troops travel along edges (large groups move slowly),
and when opposing troops meet at a node they grind each other down until one
side is gone. A player loses when they have no troops left. There is no
diplomacy and no fog of war: the entire state is visible to everyone.

\subsection{How the bots decide}
\label{sec:bg-bots}

Every bot in this report has the same skeleton, with three stages each tick:

\begin{enumerate}
  \item \textbf{Score.} A few sub-agents each score every node for how much
    they would like troops there: the \emph{economy} agent scores nodes it
    plans to build on, the \emph{expansion} agent scores capturable border
    nodes, the \emph{war} agent scores enemy-adjacent nodes. The scores are
    blended into a single target: ``here is the fraction of my army I want
    at each node.''
  \item \textbf{Move.} A transport solver turns the gap between where troops
    are and where the target wants them into actual move orders. Older bots
    nudge troops downhill along a smoothed version of this gap
    (\emph{gradient transport}); newer bots solve a routing problem that
    assigns each surplus troop a destination and a cheapest path
    (\emph{OT transport}, for optimal transport).
  \item \textbf{Act directly.} Some decisions bypass the target-and-move
    machinery: the economy agent places a building the moment a node can
    afford it, and the war agent launches an attack the moment the
    neighborhood can overwhelm a target.
\end{enumerate}

Two terms recur constantly, so we fix them here:

\begin{description}
  \item[Request.] Under OT transport, each node asks for a specific number
    of troops (``I need 40 more''), and the solver routes troops to fill
    requests. The target distribution of stage 1 is expressed through these
    requests.
  \item[Border reserve.] An always-on request at nodes that touch enemy
    territory: keep enough troops here to outnumber the largest adjacent
    enemy force. Informally: park a standing army on the border. (The code
    calls this the \emph{frontline garrison}.)
\end{description}

\subsection{The bot lineage}
\label{sec:bg-lineage}

Thirteen versions compete. Each is one deliberate change over its
predecessor; reading this table is enough to follow every result in the
report.

\begin{table}[ht]
\centering
\begin{tabular}{ll}
\toprule
Model & What changed \\
\midrule
\code{v0\_expansion}      & Baseline: economy + expansion only. Never attacks. \\
\code{v1\_knapsack}       & Adds a war agent (and drops expansion): attack \\
                          & whenever value-per-cost looks good. \\
\code{v1\_knapsack\_hybrid} & All three agents, expansion at half weight. \\
\code{v4}                 & Simpler, better war agent; the long-time default. \\
\code{v5\_qubo}           & Build layout chosen by an annealed solver. \\
\code{v6}                 & Same, with the layout biased toward factories. \\
\code{v7}                 & Stops expanding while behind on construction. \\
\code{v8}                 & Softmax over node scores becomes \emph{local}: \\
                          & troops compete for nearby targets, not global ones. \\
\code{v9}                 & Switches gradient transport for OT transport. \\
\code{v10}                & = v9 + v7's construction discipline. \\
\code{v11}                & Requests subtract troops already \emph{en route}: \\
                          & a node stops asking once help is on the way. \\
\code{v12}                & = v11 + the border reserve. \\
\bottomrule
\end{tabular}
\caption{One line per model. (\code{v2}/\code{v3} are code-identical to
\code{v4} and excluded; \code{v5\_qubo\_composite} is a \code{v5} variant
with a blended layout objective.)}
\label{tab:lineage}
\end{table}

\subsection{How we rate bots}
\label{sec:bg-rating}

Bots play round-robin tournaments (fixed seeds, both sides of every
matchup). To turn game results into a ranking we fit a Bradley--Terry model
\citep{bradleyterry1952}: assume each bot $i$ has a strength $p_i$ and that
$i$ beats $j$ with probability $p_i/(p_i+p_j)$, then find the strengths that
make the observed results most likely \citep{hunter2004}. We report
strengths on the familiar Elo scale \citep{elo1978}. This differs from
chess-style Elo updates only in being fit to the whole tournament at once:
the result does not depend on game order and is exactly reproducible from
the game log.

A one-number rating assumes strength is one-dimensional --- that if A
usually beats B and B usually beats C, then A usually beats C. That
assumption can fail (rock--paper--scissors), and \S\ref{sec:transitivity}
measures \emph{how much} it fails here: the tournament's pairwise results
are split into a part explainable by some ranking and a leftover made of
cycles (the Hodge decomposition of \citealp{jiang2011}), and the cycles are
then mapped out \citep{balduzzi2019}. The mathematics lives in
\S\ref{sec:transitivity}; nothing there is needed to read the ladder.

\subsection{Instruments}
\label{sec:bg-instruments}

Three kinds of evidence appear throughout:

\begin{itemize}
  \item \textbf{Tournaments} rate whole bots (\S\ref{sec:elo},
    \S\ref{sec:transitivity}).
  \item \textbf{Gyms} score a single component --- economy, transport,
    ordering, combat --- in isolation, on synthetic instances
    (\S\ref{sec:components}). Gym scores are diagnostics: twice in this
    report a gym and the ladder disagree, and both disagreements are
    informative.
  \item \textbf{Replays} re-run a specific game from its seed, optionally
    with one bot swapped, to explain \emph{why} a number is what it is.
    Because games are deterministic given a seed, a replay with one bot
    changed is a controlled experiment.
\end{itemize}

\paragraph{Reproducibility.} Experiments live in
\code{experiments/<name>/run.py} (fixed seeds and solver lists; run as
\code{python3 -m experiments.<name>.run}). Each writes
\code{results/<name>/}: a \code{manifest.json} recording the git SHA, date,
command, and parameters; raw per-game CSVs; and the summary tables and
figures this document includes directly. To reproduce any number: check out
the SHA in the manifest, build, re-run the module. This round is a quick
pass --- game counts give stable rankings in minutes, not tight confidence
intervals.
